% Compile with: latexmk -pdf main
\documentclass[11pt]{article}

\usepackage[margin=1in]{geometry}
\usepackage{amsmath,amssymb,amsthm}
\usepackage[T1]{fontenc}
\usepackage[utf8]{inputenc}
\usepackage{lmodern}
\usepackage{newunicodechar}
\usepackage{fancyvrb}
\usepackage{xcolor}
\usepackage[hidelinks]{hyperref}

% Unicode characters appearing in Lean code
\newunicodechar{‖}{\ensuremath{\|}}
\newunicodechar{•}{\ensuremath{\bullet}}
\newunicodechar{ℝ}{\ensuremath{\mathbb{R}}}
\newunicodechar{→}{\ensuremath{\to}}
\newunicodechar{∀}{\ensuremath{\forall}}
\newunicodechar{≠}{\ensuremath{\neq}}
\newunicodechar{≤}{\ensuremath{\leq}}
\newunicodechar{¬}{\ensuremath{\neg}}
\newunicodechar{Δ}{\ensuremath{\Delta}}
\newunicodechar{ᶠ}{\ensuremath{{}^{f}}}
\newunicodechar{𝓝}{\ensuremath{\mathcal{N}}}
\newunicodechar{²}{\ensuremath{{}^2}}
\newunicodechar{⟪}{\ensuremath{\langle\!\langle}}
\newunicodechar{⟫}{\ensuremath{\rangle\!\rangle}}

\DefineVerbatimEnvironment{lean}{Verbatim}{fontsize=\small,frame=leftline,framerule=0.6pt,
  rulecolor=\color{gray!60},xleftmargin=1em}

\newtheorem{theorem}{Theorem}
\newtheorem{lemma}[theorem]{Lemma}
\theoremstyle{remark}
\newtheorem{remark}[theorem]{Remark}

\title{A Formal Proof of Earnshaw's Theorem in Lean~4}
\author{Nichlas Löfdahl \and Claude (Anthropic)\thanks{AI system (model Claude Opus 5.5), developed by Anthropic.
Claude designed the proof strategy, wrote the Lean proofs and drafted this manuscript in collaboration with
the first author.}}
\date{October 2026}

\begin{document}
\sloppy
\maketitle

\begin{abstract}
Earnshaw's theorem states that no configuration of static charges can hold a charged particle in
stable equilibrium. Mathematically, a harmonic function has no strict local minimum or maximum. We
give a machine-checked proof in Lean~4 with Mathlib, for functions on an arbitrary
finite-dimensional real inner product space of positive dimension that are harmonic in a
neighbourhood of the point in question. Because Mathlib does not yet contain the mean-value
property or the maximum principle for harmonic functions on $\mathbb{R}^n$, we use a short
perturbation argument whose only analytic input is the second-order necessary condition for a
local minimum, which we derive from Mathlib's one-variable second-derivative test. We also prove
the weaker linear-stability statement: the Hessian of a function with vanishing Laplacian is never
positive definite. The development is about 250 lines and uses no axioms beyond Lean's standard
three.
\end{abstract}

\section{Introduction}

In 1842 Samuel Earnshaw observed that a collection of point charges interacting through
inverse-square forces cannot be in stable equilibrium~\cite{earnshaw1842}. In modern language: in
a charge-free region the electrostatic potential $V$ satisfies Laplace's equation $\Delta V=0$, and
a harmonic function has no strict local minimum. Hence the potential energy $qV$ of a test charge
has no point of stable equilibrium, whatever the sign of $q$. The theorem explains why ion traps
must use time-dependent fields (Paul traps) or magnetic fields (Penning traps), and why stable
magnetic levitation of ordinary permanent magnets is impossible without diamagnetism or feedback.

Lean's library Mathlib~\cite{mathlib2020} recently gained a definition of the Laplacian on
finite-dimensional real inner product spaces and of harmonic functions
(\texttt{InnerProductSpace.HarmonicAt}). To our knowledge, and after searching Mathlib, Physlib
and the web, Earnshaw's theorem had not been formalized in Lean or in other major proof assistants.

\paragraph{Contributions.}
\begin{itemize}
  \item A lemma identifying the second derivative of $t\mapsto g(z+tv)$ at $0$ with the Hessian
  $D^2g(z)(v,v)$, and from it the second-order necessary condition at a local minimum
  (\S\ref{sec:second}).
  \item The Laplacian of $y\mapsto\|y-x\|^2$, equal to $2\dim E$ (\S\ref{sec:quadratic}).
  \item Linear instability: if $\Delta f(x)=0$, the Hessian of $f$ at $x$ is not positive definite
  (\S\ref{sec:linear}).
  \item Earnshaw's theorem: a function harmonic near $x$ has neither a strict local minimum nor a
  strict local maximum at $x$, together with the electrostatic corollary (\S\ref{sec:earnshaw}).
\end{itemize}

The source is in the \texttt{earnshaw} folder of
\url{https://github.com/Crotalus/ai-formalized}. It builds with Lean~\texttt{v4.35.0} and
Mathlib commit \texttt{56a0058}.

\section{The argument}

Throughout, $E$ is a finite-dimensional real inner product space with $\dim E=n\ge1$.

\begin{theorem}[Earnshaw]\label{thm:earnshaw}
Let $f\colon E\to\mathbb{R}$ be $C^2$ with $\Delta f=0$ on a neighbourhood of $x$. Then $x$ is
not a strict local minimum of $f$.
\end{theorem}

\begin{proof}[Proof sketch]
Suppose $f(y)>f(x)$ for all $y\ne x$ near $x$. Choose $r>0$ such that $f$ is harmonic on the closed
ball $\bar B=\bar B(x,r)$ and $f>f(x)$ on $\bar B\setminus\{x\}$. The sphere $S=\partial B(x,r)$
is compact and non-empty, so $f$ attains a minimum $f(y_0)$ on $S$; put $\delta=f(y_0)-f(x)>0$.
Let
\[
  g(y)=f(y)-\varepsilon\|y-x\|^2,\qquad \varepsilon=\frac{\delta}{2r^2}.
\]
Then $g(x)=f(x)$, while on $S$ we have $g\ge f(y_0)-\delta/2=f(x)+\delta/2$. Since $\bar B$ is
compact, $g$ attains its minimum on $\bar B$ at some $z$, with $g(z)\le g(x)<\min_S g$, so $z$ lies
in the open ball and is a local minimum of $g$. By the second-order necessary condition
(Lemma~\ref{lem:soc}), $D^2g(z)(v,v)\ge0$ for all $v$, and summing over an orthonormal basis gives
$\Delta g(z)\ge0$. But
$\Delta g(z)=\Delta f(z)-\varepsilon\,\Delta\|\cdot-x\|^2(z)=0-2n\varepsilon<0$, a contradiction.
\end{proof}

The maximum version follows by applying the theorem to $-f$. This argument is a standard device in
the theory of maximum principles (cf.\ the proof of the weak maximum principle for elliptic operators in~\cite{evans}). Its appeal for formalization
is that it needs neither integration over spheres nor the mean-value property.

\begin{remark}
The hypothesis $\dim E\ge1$ is necessary: if $E=\{0\}$, then $x$ has no punctured neighbourhood,
so ``strict local minimum'' holds vacuously. In Lean this appears as \texttt{[Nontrivial E]}.
\end{remark}

\section{Formalization}

\subsection{Second derivatives along lines}\label{sec:second}

Mathlib defines the Laplacian through the second iterated Fréchet derivative and proves
\[
  \Delta f(x)=\sum_i D^2f(x)(e_i,e_i)
\]
for any orthonormal basis $(e_i)$ (\texttt{laplacian\_eq\_iteratedFDeriv\_orthonormalBasis}). The
missing ingredient is the link between $D^2g(z)(v,v)$ and one-variable calculus:

\begin{lean}
lemma deriv_deriv_comp_line {g : E → ℝ} {z : E} (hg : ContDiffAt ℝ 2 g z) (v : E) :
    deriv (deriv fun t : ℝ => g (z + t • v)) 0 = iteratedFDeriv ℝ 2 g z ![v, v]
\end{lean}

The proof first shows that $\frac{d}{dt}g(z+tv)=Dg(z+tv)\,v$ for $t$ near $0$. This holds because
$g$ is differentiable on a neighbourhood of $z$, which follows from \texttt{ContDiffAt.eventually}.
It then differentiates once more, using that $Dg$ is differentiable at $z$
(\texttt{ContDiffAt.fderiv\_right}) and the product rule for applying a continuous linear map.

\begin{lemma}[Second-order necessary condition]\label{lem:soc}
If $g$ is $C^2$ at $z$ and $z$ is a local minimum of $g$, then $D^2g(z)(v,v)\ge0$ for all $v$.
\end{lemma}

Mathlib has the \emph{sufficient} second-derivative test for functions $\mathbb{R}\to\mathbb{R}$
(\texttt{isLocalMax\_of\_deriv\_deriv\_neg}), but not the necessary condition. We derive the latter
with a small trick. Let $\varphi(t)=g(z+tv)$, which has a local minimum at $0$, so $\varphi'(0)=0$.
If $\varphi''(0)<0$, the sufficient test makes $0$ also a local maximum. Then $\varphi$ is
constant near $0$, so $\varphi'$ vanishes near $0$ and $\varphi''(0)=0$, a contradiction. In Lean
this is about fifteen lines.

\subsection{The squared distance}\label{sec:quadratic}

\begin{lean}
lemma laplacian_norm_sub_sq (x z : E) :
    Δ (fun y => ‖y - x‖ ^ 2) z = 2 * Module.finrank ℝ E
\end{lean}

Rather than computing the Hessian of $y\mapsto\|y-x\|^2$ abstractly, we reuse
\texttt{deriv\_deriv\_comp\_line}. Along a line,
$\|z+tv-x\|^2=\|z-x\|^2+2t\langle z-x,v\rangle+t^2\|v\|^2$, a quadratic polynomial in $t$ with
second derivative $2\|v\|^2$. Summing over an orthonormal basis gives $2n$.

\subsection{Linear instability}\label{sec:linear}

\begin{lean}
lemma not_posDef_hessian_of_laplacian_eq_zero [Nontrivial E] {f : E → ℝ} {x : E}
    (h : Δ f x = 0) : ¬ ∀ v ≠ 0, 0 < iteratedFDeriv ℝ 2 f x ![v, v]
\end{lean}

If the Hessian were positive definite, the Laplacian would be a non-empty sum of positive terms.
This is the version of Earnshaw's theorem usually given in physics texts. It rules out equilibria
that are stable to second order, but on its own it does not exclude a degenerate strict minimum.
For example, $t\mapsto t^4$ has a strict minimum at $0$ with vanishing second derivative.

\subsection{Earnshaw's theorem}\label{sec:earnshaw}

\begin{lean}
theorem HarmonicAt.not_isStrictLocalMin [Nontrivial E] {f : E → ℝ} {x : E}
    (hf : HarmonicAt f x) : ¬ ∀ᶠ y in 𝓝[≠] x, f x < f y

theorem HarmonicAt.not_isStrictLocalMax [Nontrivial E] {f : E → ℝ} {x : E}
    (hf : HarmonicAt f x) : ¬ ∀ᶠ y in 𝓝[≠] x, f y < f x

theorem earnshaw {V : EuclideanSpace ℝ (Fin 3) → ℝ} {x : EuclideanSpace ℝ (Fin 3)}
    (hV : HarmonicAt V x) (q : ℝ) : ¬ ∀ᶠ y in 𝓝[≠] x, q * V x < q * V y
\end{lean}

Mathlib's \texttt{HarmonicAt f x} means that $f$ is $C^2$ at $x$ and $\Delta f=0$ on a
neighbourhood of $x$. It is an open condition (\texttt{HarmonicAt.eventually}), which supplies the
radius on which $f$ is harmonic. The formal proof follows the sketch above step by step. It
extracts radii from the two neighbourhood filters, takes minima over the compact sphere and closed
ball (\texttt{IsCompact.exists\_isMinOn}), shows that the minimiser of $g$ is not on the sphere,
and closes with the linearity of the Laplacian (\texttt{ContDiffAt.laplacian\_sub},
\texttt{laplacian\_smul}) and \texttt{nlinarith}. The main proof is 47 lines.

The physical corollary covers both signs of the test charge at once, because
\texttt{HarmonicAt} is preserved under scalar multiplication.

\section{Discussion}

\paragraph{Size and checking.} The development has 9 lemmas and theorems in about 250 lines,
including documentation. It compiles without warnings, and \texttt{\#print axioms} on the main
results reports only \texttt{propext}, \texttt{Classical.choice} and \texttt{Quot.sound}.

\paragraph{What made it feasible.} Three pieces of Mathlib carried the weight: the Laplacian and
its expression in an orthonormal basis, the notion of a harmonic function with its stability under
negation and scaling, and the one-variable second-derivative test. The two lemmas we contribute,
the Hessian along a line and the necessary condition at a minimum, are general facts that would
fit naturally into Mathlib.

\paragraph{Limitations and future work.} We do not prove the mean-value property or the strong
maximum principle, which would yield stronger statements (for example, that a harmonic function
with a non-strict local minimum is locally constant). The electrostatic corollary assumes that the
potential is harmonic. Deriving this from Coulomb's law for a finite charge distribution away from
the charges would connect the result to Physlib's electromagnetism library. Natural extensions are
the magnetostatic version (Braunbek's theorem on diamagnetic levitation) and a dynamical
statement: Lyapunov instability of the equilibrium for the Newtonian motion $m\ddot{\mathbf r}=-q\nabla V$.

\paragraph{Author contributions.} N.L.\ chose the problem, directed the work and reviewed the
results. Claude, an AI system developed by Anthropic, surveyed prior work, devised the
perturbation strategy and the second-order-condition trick, wrote the Lean proofs and drafted the
manuscript. All statements were checked by the Lean kernel. As an AI system cannot take
responsibility for a publication, N.L.\ is accountable for its content.

\bibliographystyle{plain}
\bibliography{refs}

\end{document}
